\subsubsection{Penalty Coefficients}\label{sec:penalty-coefficients}

Suppose our original objective is simply:
\begin{equation*}
    \min y(x)
\end{equation*}
And we have a constraint:
\begin{equation*}
    x_1 + x_2 \le 1
\end{equation*}
We know the penalty is (\autopageref{ex:penalty-terms-x1+x2-le-1}):
\begin{equation*}
    \text{penalty}(x) = x_1 x_2
\end{equation*}
However, in QUBO we need to minimize:
\begin{equation*}
    \min y(x) + \gamma \cdot \text{penalty}(x)
\end{equation*}
Where $\gamma$ is the \textbf{penalty coefficient}.

\highspace
\textcolor{Green3}{\faIcon{question-circle} \textbf{Why do we need $\gamma$?}} Because the objective function and the penalty may have \textbf{completely different scales}. For example, let's say that the feasible solution has $y(x) = 100$ and a penalty of $0$, while the infeasible solution has a much smaller objective value, say $y(x) = 1$, but a penalty of $1$. \textbf{Without} a cofficient, the \hl{infeasible solution would be preferred}, which is not what we want.

\highspace
\begin{flushleft}
    \textcolor{Red2}{\faIcon{exclamation-triangle} \textbf{$\gamma$ must be chosen carefully!}}
\end{flushleft}
\begin{itemize}
    \item Penalty \textbf{too small}. Suppose we have $\gamma = 0.1$. Then, the infeasible solution would have a total cost of $1 + 0.1 \cdot 1 = 1.1$, which is still less than the feasible solution's cost of $100$. Therefore, the infeasible solution would still be preferred. In general, \hl{if the penalty is too small, it may lead to the selection of unfeasible solutions}.


    \item Penalty \textbf{large enough}. Suppose we have $\gamma = 100$. Then, the infeasible solution would have a total cost of $1 + 100 \cdot 1 = 101$, which is greater than the feasible solution's cost of $100$. Therefore, the feasible solution would be preferred. In general, \hl{if the penalty is large enough, it will ensure that feasible solutions are preferred over infeasible ones}.


    \item Penalty \textbf{too large}. Suppose we have $\gamma = 1000$. Then, the infeasible solution would have a total cost of $1 + 1000 \cdot 1 = 1001$, which is much greater than the feasible solution's cost of $100$. Therefore, the feasible solution would be strongly preferred. In general, \hl{if the penalty is too large, it may overwhelm the problem and lead to non optimal solutions}.
\end{itemize}
In practice, $\gamma$ is simply an \textbf{hyperparameter} because there is no formula that tells us the perfect value for it. The modeler chooses it. However, we can \textbf{start with a penalty coefficient close to one and fine tune from there}.